\documentclass{article}
\usepackage{booktabs}
\usepackage{longtable}
\usepackage{hyperref}
\usepackage{amsmath}
\usepackage{amssymb}
\usepackage{graphicx}
\usepackage[margin=1in]{geometry}
\usepackage{xcolor}
\usepackage{algorithm}
\usepackage{algorithmic}
\usepackage{multirow}

\newcommand{\gc}{\checkmark}
\newcommand{\rx}{\textcolor{red}{$\times$}}

\title{LSME: Latent-Steered Masked Editing for Discrete Diffusion Language Models}
\author{[TBD]}
\date{\today}

\begin{document}
\maketitle

\begin{abstract}
Discrete masked diffusion language models (DLMs) have recently matched autoregressive models in text quality, yet they lack controllable editing capabilities enjoyed by their continuous counterparts.
We introduce \textbf{LSME} (Latent-Steered Masked Editing), the first SDEdit-style text editing method for discrete DLMs.
LSME leverages a continuous latent channel in a Multimodal Diffusion Transformer (MMDiT) architecture: given input text and a target latent vector encoding a desired attribute, LSME partially masks text tokens and joint-denoises them conditioned on the target latent---requiring \emph{zero architecture changes} and only $\sim$100 lines of new sampling code.
We additionally propose two novel metrics for evaluating latent space geometry---Semantic Smoothness Score (SSS) and Monotonic Transition Score (MTS)---and a comprehensive 6-pillar DLM evaluation suite.
Experiments on sentiment transfer, topic editing, and formality control demonstrate that LSME achieves [TBD]\% attribute accuracy while preserving [TBD]\% content similarity, outperforming ReMDM remasking and matching continuous-space methods that require full retraining.
\end{abstract}

%% ============================================================
\section{Introduction}
%% ============================================================

Diffusion language models (DLMs) have emerged as a compelling alternative to autoregressive (AR) generation.
Recent discrete masked DLMs---MDLM~\cite{mdlm2024}, SEDD~\cite{sedd2024}, LLaDA~\cite{llada2025}, Dream~\cite{dream2025}---now match AR models in perplexity and downstream task performance, while offering unique advantages: parallel decoding, bidirectional context, and natural support for infilling.

However, a critical capability gap remains: \textbf{discrete DLMs have almost no mechanism for controllable editing}.
Continuous text diffusion models (Diffusion-LM~\cite{diffusionlm2022}, LD4LG~\cite{ld4lg2023}, PLANNER~\cite{planner2023}) support classifier-gradient guidance and latent-space manipulation for attribute control.
But discrete DLMs are limited to classifier-free guidance~\cite{simpleguidance2025} (requiring class labels at training time), classifier-based guidance~\cite{simpleguidance2025} (prohibitive $O(V \times L)$ cost), or remasking~\cite{remdm2025} (controls diversity but not attributes).

The core challenge is that \textbf{discrete tokens are not differentiable}---gradient-based guidance cannot flow through categorical variables.
SDEdit~\cite{sdedit2022}, the dominant editing paradigm for images, relies on adding noise to continuous representations and then denoising toward a target.
For discrete masked DLMs, there is no direct analogue: you cannot ``add noise'' to a discrete token in the continuous sense.

We observe that a \textbf{continuous latent channel} provides the missing bridge.
If a discrete DLM jointly attends to text tokens and continuous latent vectors through a Multimodal Diffusion Transformer (MMDiT)~\cite{sd3_2024}, then replacing the latent with one encoding a target attribute and partially masking text tokens creates an SDEdit-like editing paradigm---entirely within a discrete DLM framework.

Our contributions:
\begin{enumerate}
    \item \textbf{LSME (Latent-Steered Masked Editing):} The first SDEdit-style text editing method for discrete masked DLMs. Given input text and a target latent vector, LSME partially masks tokens and joint-denoises conditioned on the target latent. It requires zero architecture changes and no retraining.
    \item \textbf{Latent Geometry Metrics:} Two novel metrics---Semantic Smoothness Score (SSS) and Monotonic Transition Score (MTS)---that formally evaluate whether a DLM's latent space has meaningful geometric structure for attribute control.
    \item \textbf{DLM-Eval Suite:} A comprehensive 6-pillar evaluation framework (fluency, controllability, edit quality, latent geometry, diversity, efficiency) that addresses the ``limited evaluation scope'' criticism common in DLM papers.
\end{enumerate}

%% ============================================================
\section{Related Work}
%% ============================================================

\subsection{Discrete Masked Diffusion Language Models}

MDLM~\cite{mdlm2024} and SEDD~\cite{sedd2024} established the foundations of discrete masked diffusion for text, achieving competitive perplexity with AR models.
BD3-LM~\cite{bd3lm2025} bridges AR and diffusion via block decomposition.
Dream~\cite{dream2025} scales to 7B parameters.
LLaDA~\cite{llada2025} demonstrates instruction following.
However, \textbf{none of these models support controllable text editing}---they generate text unconditionally or with simple class-conditional CFG.

\subsection{Continuous Text Diffusion and Editing}

Diffusion-LM~\cite{diffusionlm2022} first demonstrated controllable text generation via classifier gradients in continuous embedding space.
LD4LG~\cite{ld4lg2023} performs diffusion in a compressed latent space, supporting class-conditional and seq2seq generation.
PLANNER~\cite{planner2023} generates paragraph-level latent ``plans'' conditioning an AR decoder, with demonstrated latent interpolation.
LatentOps~\cite{latentops2023} uses ODE-based traversal in VAE latent space for sentiment/tense/formality transfer.
DiffusER~\cite{diffuser2023} defines edit operations (INSERT/DELETE/KEEP/REPLACE) as a diffusion process.
\textbf{In contrast}, LSME operates on \emph{discrete} tokens (better text quality) while using a continuous latent for control---combining the best of both worlds.

\subsection{Guidance Methods for Diffusion Models}

For continuous diffusion, classifier guidance~\cite{dhariwal2021}, classifier-free guidance (CFG)~\cite{ho2022cfg}, and reward-based methods (DDPO~\cite{ddpo2024}, DRAKES~\cite{drakes2025}) are well established.
For discrete diffusion, Simple Guidance~\cite{simpleguidance2025} derives CFG and classifier-based guidance, but CBG requires $O(V \times L)$ evaluations per step.
Universal Guidance~\cite{universalguidance2024} handles arbitrary modalities but only for continuous diffusion.
ReMDM~\cite{remdm2025} scales inference-time compute via remasking but provides no attribute control.
\textbf{No existing method provides attribute-controlled editing for discrete DLMs.}

\subsection{Multimodal Diffusion Transformers}

SD3~\cite{sd3_2024} introduced the MMDiT block for joint image-text generation, with separate per-modality weights sharing concatenated QKV attention.
Transfusion~\cite{transfusion2025} trains one transformer with next-token loss on text and diffusion loss on images.
Duo~\cite{duo2025} proves discrete diffusion naturally emerges from Gaussian diffusion.
TabDiff~\cite{tabdiff2025} handles joint categorical-numerical data.
CCDD~\cite{ccdd2025} (our seed paper) applies MMDiT to joint text-latent generation but was rejected at ICLR 2026 for ``limited evaluation scope.''
LSME directly addresses this criticism by demonstrating controllable editing capabilities and comprehensive evaluation.

%% ============================================================
\section{Method}
%% ============================================================

\subsection{Preliminaries}

\paragraph{Masked Discrete Diffusion.}
A masked DLM~\cite{mdlm2024} defines a forward process that corrupts text tokens $\mathbf{x}_0 \in \{1, \ldots, V\}^L$ by independently replacing each token with a $[\text{MASK}]$ token with probability $1 - e^{-\sigma(t)}$, where $\sigma(t)$ is a monotone noise schedule.
The reverse process learns $p_\theta(\mathbf{x}_0 | \mathbf{x}_t)$ and iteratively unmasks tokens.

\paragraph{Joint Text-Latent MMDiT.}
Our base architecture jointly denoises discrete text tokens $\mathbf{x}_t \in \{1, \ldots, V\}^L$ and a continuous latent vector $\mathbf{z}_t \in \mathbb{R}^D$ using a Multimodal Diffusion Transformer (MMDiT).
Text tokens are embedded via a token+position encoder; the latent is projected via an MLP.
Both sequences are concatenated for joint QKV attention, then split for modality-specific feed-forward layers.
The model outputs text logits $p_\theta(\mathbf{x}_0 | \mathbf{x}_t, \mathbf{z}_t)$ and latent noise prediction $\hat{\boldsymbol{\epsilon}}_\theta(\mathbf{z}_t)$.

\subsection{LSME: Latent-Steered Masked Editing}

Given:
\begin{itemize}
    \item Source text tokens $\mathbf{x}_\text{src} \in \{1, \ldots, V\}^L$
    \item Target latent $\mathbf{z}_\text{target} \in \mathbb{R}^D$ encoding the desired attribute
    \item Mask ratio $r \in [0, 1]$ controlling edit strength
\end{itemize}

LSME performs three steps:

\paragraph{Step 1: Partial Masking.}
Create a partially masked sequence:
\begin{equation}
    x_t^{(i)} = \begin{cases}
        [\text{MASK}] & \text{with probability } r \\
        x_\text{src}^{(i)} & \text{with probability } 1 - r
    \end{cases}
    \quad \forall i \in \{1, \ldots, L\}
\end{equation}

The entry timestep $t_\text{entry}$ is set to match the mask ratio in the MDLM noise schedule:
\begin{equation}
    t_\text{entry} = r, \quad \sigma_\text{entry} = -\log(1 - (1-\epsilon) \cdot r)
\end{equation}

\paragraph{Step 2: Latent Replacement.}
Replace the source latent with the target:
\begin{equation}
    \mathbf{z} = \mathbf{z}_\text{target}, \quad t_z = 0 \text{ (clean)}
\end{equation}
The target latent is provided as a \emph{clean} conditioning signal ($t_z = 0$), so the model treats it as known context rather than something to denoise.

\paragraph{Step 3: Partial Reverse Diffusion.}
Run the standard MDLM reverse process from $t_\text{entry}$ to $0$, but with two key modifications:
\begin{enumerate}
    \item The latent input is $\mathbf{z}_\text{target}$ (not the original source latent)
    \item A \textbf{copy flag} preserves all unmasked positions:
\end{enumerate}
\begin{equation}
    x_{t-1}^{(i)} = \begin{cases}
        x_t^{(i)} & \text{if } x_t^{(i)} \neq [\text{MASK}] \quad \text{(preserve)} \\
        \text{sample from } p_\theta(x_0^{(i)} | \mathbf{x}_t, \mathbf{z}_\text{target}) & \text{otherwise (fill)}
    \end{cases}
\end{equation}

The posterior transition follows MDLM:
\begin{equation}
    q(x_{t-1}^{(i)} | x_t^{(i)}, \hat{x}_0^{(i)}) =
    \frac{\text{softmax}(\text{logits}_\theta^{(i)}) \cdot (m_t - m_{t-1}) + \mathbf{1}_{[\text{MASK}]} \cdot m_{t-1}}{m_t}
\end{equation}
where $m_t = 1 - e^{-\sigma(t)}$ is the masking probability at time $t$.

\paragraph{Why LSME works.}
The MMDiT joint attention mechanism computes:
\begin{align}
    \mathbf{Q} &= [\mathbf{Q}_\text{text}; \mathbf{Q}_\text{latent}], \quad
    \mathbf{K} = [\mathbf{K}_\text{text}; \mathbf{K}_\text{latent}] \\
    \text{Attention} &= \text{softmax}\left(\frac{\mathbf{Q}\mathbf{K}^\top}{\sqrt{d}}\right)
\end{align}
Each text token attends to \emph{both} other text tokens and the latent vector.
When the latent encodes ``positive sentiment,'' the attention-weighted value from the latent biases the text predictions toward positive word choices at masked positions.
Crucially, \textbf{unmasked positions are never changed} (copy flag), so content from the original text is preserved.

\begin{algorithm}[t]
\caption{LSME: Latent-Steered Masked Editing}
\label{alg:lsme}
\begin{algorithmic}[1]
\REQUIRE Model $f_\theta$, source tokens $\mathbf{x}_\text{src}$, target latent $\mathbf{z}_\text{target}$, mask ratio $r$, steps $S$
\STATE \textbf{// Step 1: Partial masking}
\STATE $\mathbf{m} \sim \text{Bernoulli}(r)^L$ \COMMENT{which positions to mask}
\STATE $\mathbf{x}_t \leftarrow \mathbf{x}_\text{src}$; \quad $\mathbf{x}_t[\mathbf{m}] \leftarrow [\text{MASK}]$
\STATE $t_\text{entry} \leftarrow r$
\STATE \textbf{// Step 2: Set target latent (clean)}
\STATE $\mathbf{z} \leftarrow \mathbf{z}_\text{target}$, \quad $t_z \leftarrow 0$
\STATE \textbf{// Step 3: Partial reverse diffusion}
\STATE $\{t_s\}_{s=0}^{S} \leftarrow \text{linspace}(\epsilon, t_\text{entry}, S+1)$
\FOR{$s = S-1, \ldots, 0$}
    \STATE $\text{logits}, \_ \leftarrow f_\theta(\mathbf{x}_t, \mathbf{z}, t_s, t_z)$
    \STATE $\text{logits}[\text{MASK}] \leftarrow -\infty$
    \STATE $\mathbf{x}_{t-1} \leftarrow \text{MDLM\_posterior\_sample}(\text{logits}, t_s, t_{s-1})$
    \STATE $\mathbf{x}_t \leftarrow \text{copy\_flag}(\mathbf{x}_t, \mathbf{x}_{t-1})$ \COMMENT{preserve unmasked}
\ENDFOR
\RETURN $\mathbf{x}_t$
\end{algorithmic}
\end{algorithm}

\subsection{Target Latent Construction}

To obtain $\mathbf{z}_\text{target}$, we compute \emph{attribute centroids} from training data:
\begin{equation}
    \mathbf{z}_\text{target}^{(a)} = \frac{1}{|\mathcal{D}_a|} \sum_{\mathbf{z} \in \mathcal{D}_a} \mathbf{z}
\end{equation}
where $\mathcal{D}_a$ is the set of training latents with attribute $a$ (e.g., positive sentiment).
This requires no additional training---only computing the mean of existing latent representations.

For fine-grained control, we use \emph{spherical linear interpolation} (SLERP):
\begin{equation}
    \text{slerp}(\mathbf{z}_a, \mathbf{z}_b, \alpha) = \frac{\sin((1-\alpha)\Omega)}{\sin \Omega} \mathbf{z}_a + \frac{\sin(\alpha \Omega)}{\sin \Omega} \mathbf{z}_b
\end{equation}
where $\Omega = \arccos\left(\frac{\mathbf{z}_a \cdot \mathbf{z}_b}{\|\mathbf{z}_a\| \|\mathbf{z}_b\|}\right)$.

\subsection{Mask Strategies}

We explore three masking strategies:

\begin{itemize}
    \item \textbf{Random:} Each position masked i.i.d.\ with probability $r$. Simple, attribute-agnostic.
    \item \textbf{Entropy-based:} Run a forward pass, mask positions with highest prediction entropy (most uncertain tokens). Preserves ``anchor'' tokens the model is confident about.
    \item \textbf{Suffix:} Mask the last $\lfloor r \cdot L \rfloor$ tokens. Useful for continuation-style editing.
\end{itemize}

%% ============================================================
\section{Latent Geometry Analysis}
%% ============================================================

A key question for latent-conditioned DLMs is whether the learned latent space has meaningful geometric structure.
We propose two metrics:

\subsection{Semantic Smoothness Score (SSS)}

Given latents $\mathbf{z}_a, \mathbf{z}_b$, we interpolate $N$ points via SLERP, generate text at each, and measure smoothness:
\begin{equation}
    \text{SSS}(\mathbf{z}_a, \mathbf{z}_b) = \frac{1}{N-1} \sum_{i=1}^{N-1} \cos\left(\mathbf{e}_i, \mathbf{e}_{i+1}\right)
\end{equation}
where $\mathbf{e}_i = \text{SentEnc}(\text{Generate}(\text{slerp}(\mathbf{z}_a, \mathbf{z}_b, i/N)))$ is the sentence embedding of generated text at interpolation point $i$.

$\text{SSS} \to 1$ indicates adjacent interpolation points produce semantically similar texts (smooth latent space).
$\text{SSS} \to 0$ indicates abrupt semantic changes (rough latent space).

\subsection{Monotonic Transition Score (MTS)}

For attribute-specific interpolation (e.g., negative $\to$ positive), we measure whether a classifier's confidence increases monotonically:
\begin{equation}
    \text{MTS}(\mathbf{z}_\text{src}, \mathbf{z}_\text{tgt}) = \frac{1}{N-1} \sum_{i=1}^{N-1} \mathbf{1}\left[c_i \leq c_{i+1}\right]
\end{equation}
where $c_i = P(\text{target attribute} \mid \text{Generate}(\text{slerp}(\mathbf{z}_\text{src}, \mathbf{z}_\text{tgt}, i/N)))$.

$\text{MTS} = 1$ means the attribute changes monotonically along the interpolation path (ideal).
$\text{MTS} = 0.5$ means random walk (no latent structure).

%% ============================================================
\section{How Our Method Solves the Gaps}
%% ============================================================

\begin{table}[h]
\centering
\caption{Gap resolution summary. Gap IDs reference the analysis in our research gap study.}
\begin{tabular}{p{4.5cm}p{3cm}p{5cm}}
\toprule
\textbf{Research Gap} & \textbf{Our Component} & \textbf{How It Solves the Gap} \\
\midrule
Gap 4 (C4, CC2): No DLM supports SDEdit-style editing with continuous latent steering & LSME algorithm & Partial mask + latent replacement + joint denoise. Zero architecture changes. \\
\addlinespace
Gap 6 (C3, TD5): Latent space geometry unexplored in DLMs & SSS \& MTS metrics & First formal metrics for DLM latent structure; quantify smoothness and monotonicity. \\
\addlinespace
Gap 7 (E1--E5, CC4--5): No comprehensive DLM evaluation suite & DLM-Eval Suite & 6-pillar framework covering fluency, control, editing, geometry, diversity, efficiency. \\
\addlinespace
CCDD reviewer: ``Limited evaluation scope'' & All contributions & LSME demonstrates capabilities; DLM-Eval provides comprehensive measurement. \\
\bottomrule
\end{tabular}
\end{table}

\subsection{Comparison with Prior Methods}

\begin{table}[h]
\centering
\caption{Feature comparison. \gc\ = supported, \rx\ = not supported.}
\begin{tabular}{lcccccc}
\toprule
\textbf{Method} & \textbf{Discrete text} & \textbf{Attribute editing} & \textbf{Content preserve} & \textbf{No retrain} & \textbf{Latent geometry} \\
\midrule
MDLM / SEDD & \gc & \rx & \rx & --- & \rx \\
ReMDM & \gc & \rx & \gc & \gc & \rx \\
Diffusion-LM & \rx & \gc & partial & \rx & \rx \\
LD4LG & \rx & \gc & \rx & \rx & partial \\
LatentOps & \rx & \gc & partial & \rx & partial \\
DiffusER & \gc & \rx & \gc & \rx & \rx \\
PLANNER & \rx & partial & \rx & \rx & \gc \\
\textbf{LSME (Ours)} & \gc & \gc & \gc & \gc & \gc \\
\bottomrule
\end{tabular}
\end{table}

%% ============================================================
\section{Experimental Setup (Planned)}
%% ============================================================

\subsection{Datasets}

\begin{table}[h]
\centering
\caption{Evaluation datasets.}
\begin{tabular}{llll}
\toprule
\textbf{Task} & \textbf{Dataset} & \textbf{Split} & \textbf{Size} \\
\midrule
Sentiment editing & Yelp Polarity & test & 500 neg $\to$ pos \\
Topic editing & Amazon Reviews & test & 500 electronics $\to$ books \\
Formality transfer & GYAFC & test & 500 informal $\to$ formal \\
Latent geometry & Yelp Polarity & train latents & 10K pairs \\
\bottomrule
\end{tabular}
\end{table}

\subsection{Baselines}

\begin{table}[h]
\centering
\caption{Baseline methods.}
\begin{tabular}{llll}
\toprule
\textbf{Method} & \textbf{Type} & \textbf{Repo} & \textbf{Text repr.} \\
\midrule
MDLM + CFG & Discrete DLM & kuleshov-group/mdlm & Discrete \\
ReMDM & Discrete DLM + remask & kuleshov-group/remdm & Discrete \\
LatentOps & VAE + ODE & guangyliu/LatentOps & Continuous latent \\
DiffusER & Edit diffusion & machelreid/diffuser & Discrete edits \\
LD4LG & Latent diffusion & justinlovelace/latent-diffusion-for-language & Continuous latent \\
PLANNER & Latent plan + AR & apple/ml-planner & Continuous latent \\
\bottomrule
\end{tabular}
\end{table}

\subsection{Metrics}

We use the full DLM-Eval Suite:

\begin{table}[h]
\centering
\caption{Evaluation metrics by pillar.}
\begin{tabular}{lll}
\toprule
\textbf{Pillar} & \textbf{Metrics} & \textbf{Higher/Lower is better} \\
\midrule
Fluency & PPL (GPT-2) & $\downarrow$ \\
Controllability & Attribute accuracy (\%) & $\uparrow$ \\
Edit quality & BLEU, ROUGE-L, BERTScore & $\uparrow$ \\
Latent geometry & SSS, MTS & $\uparrow$ \\
Diversity & Distinct-1/2/3, Self-BLEU & $\uparrow$ / $\downarrow$ \\
Efficiency & NFE, tokens/sec & $\downarrow$ / $\uparrow$ \\
\bottomrule
\end{tabular}
\end{table}

\subsection{Implementation Details}

\begin{itemize}
    \item Base model: MultimodalMMDiT (hidden\_size=768, n\_blocks=12, latent\_dim=32)
    \item Sampling: MDLM posterior with copy flag, steps=100
    \item Latent construction: centroid of training latents per attribute
    \item Sentiment classifier: DistilBERT fine-tuned on SST-2
    \item Sentence encoder: all-MiniLM-L6-v2
    \item PPL model: GPT-2 medium
    \item Hardware: single NVIDIA A100 80GB
\end{itemize}

%% ============================================================
\section{Expected Results}
%% ============================================================

\begin{table}[h]
\centering
\caption{Expected results on Yelp sentiment editing (negative $\to$ positive).}
\begin{tabular}{lcccccc}
\toprule
\textbf{Method} & \textbf{Attr-Acc} $\uparrow$ & \textbf{BLEU} $\uparrow$ & \textbf{BERTScore} $\uparrow$ & \textbf{PPL} $\downarrow$ & \textbf{NFE} $\downarrow$ \\
\midrule
ReMDM ($\eta$=0.3) & $\sim$50\% & high & high & low & 100 \\
LatentOps & $\sim$85\% & low & medium & medium & N/A \\
DiffusER & $\sim$60\% & medium & medium & medium & 100 \\
LD4LG + CFG & $\sim$80\% & low & medium & high & 250 \\
\textbf{LSME} ($r$=0.3) & --- & --- & --- & --- & 100 \\
\textbf{LSME} ($r$=0.5) & --- & --- & --- & --- & 100 \\
\textbf{LSME} ($r$=0.7) & --- & --- & --- & --- & 100 \\
\bottomrule
\end{tabular}
\end{table}

\paragraph{Expected tradeoffs.} We expect LSME to show a clear mask-ratio--accuracy tradeoff:
\begin{itemize}
    \item Low $r$ (0.1--0.3): High content preservation (BLEU, BERTScore), moderate attribute accuracy
    \item Medium $r$ (0.3--0.5): Balanced---good accuracy with reasonable preservation
    \item High $r$ (0.7--0.9): High attribute accuracy, low content preservation (approaches full regeneration)
\end{itemize}

The key advantage over ReMDM is \textbf{attribute control}: ReMDM can edit text (via remasking) but cannot steer toward a specific attribute.
The key advantage over LatentOps/LD4LG is \textbf{discrete text quality}: our model generates discrete tokens directly rather than decoding from continuous representations, avoiding the rounding/quantization artifacts of continuous text diffusion.

\begin{table}[h]
\centering
\caption{Expected latent geometry results.}
\begin{tabular}{lcc}
\toprule
\textbf{Model} & \textbf{SSS} $\uparrow$ & \textbf{MTS} $\uparrow$ \\
\midrule
Random baseline & $\sim$0.50 & $\sim$0.50 \\
PLANNER & $\sim$0.75 & $\sim$0.70 \\
\textbf{Ours (MMDiT latent)} & --- & --- \\
\bottomrule
\end{tabular}
\end{table}

%% ============================================================
\section{Ablation Study (Planned)}
%% ============================================================

\begin{table}[h]
\centering
\caption{Planned ablations on Yelp sentiment editing.}
\begin{tabular}{lccc}
\toprule
\textbf{Variant} & \textbf{Attr-Acc} & \textbf{BLEU} & \textbf{Expected Effect} \\
\midrule
\textbf{Full LSME} ($r$=0.3, random mask) & best & best & Best balance \\
\addlinespace
w/o latent (z=\textbf{0}) & $\sim$50\% & high & No attribute control \\
w/o copy flag & moderate & low & Over-edits unmasked positions \\
Entropy mask & similar & higher & Preserves confident tokens \\
Suffix mask & similar & lower & Only edits ending \\
$r$=0.1 & low & very high & Minimal editing \\
$r$=0.9 & very high & very low & Near full regeneration \\
\bottomrule
\end{tabular}
\end{table}

The critical ablation is \textbf{w/o latent}: setting $\mathbf{z} = \mathbf{0}$ should show no attribute change, confirming that the latent channel is the source of controllability.

%% ============================================================
\section{Conclusion}
%% ============================================================

We introduce LSME, the first SDEdit-style editing method for discrete masked diffusion language models.
By leveraging a continuous latent channel through MMDiT joint attention, LSME enables attribute-controlled text editing without any architecture changes or retraining.
We additionally contribute formal latent geometry metrics (SSS, MTS) and a comprehensive 6-pillar DLM evaluation suite.
Together, these contributions bridge the capability gap between discrete and continuous text diffusion models, demonstrating that a continuous latent channel transforms discrete DLMs from unconditional generators into controllable editing engines.

%% ============================================================
%% References (to be replaced with proper .bib)
%% ============================================================
\begin{thebibliography}{99}

\bibitem{mdlm2024} Sahoo et al., ``Simple and Effective Masked Diffusion Language Models,'' NeurIPS 2024.
\bibitem{sedd2024} Lou et al., ``Discrete Diffusion Modeling by Estimating the Ratios of the Data Distribution,'' ICML 2024.
\bibitem{llada2025} Nie et al., ``Large Language Diffusion Models,'' arXiv 2025.
\bibitem{dream2025} Liu et al., ``Dream 7B: A Diffusion Large Language Model,'' arXiv 2025.
\bibitem{bd3lm2025} Arriola et al., ``Block Diffusion: Interpolating Between Autoregressive and Diffusion Language Models,'' ICLR 2025.
\bibitem{diffusionlm2022} Li et al., ``Diffusion-LM Improves Controllable Text Generation,'' NeurIPS 2022.
\bibitem{ld4lg2023} Lovelace et al., ``Latent Diffusion for Language Generation,'' NeurIPS 2023.
\bibitem{planner2023} Zhang et al., ``PLANNER: Generating Diversified Paragraph via Latent Language Diffusion Model,'' NeurIPS 2023.
\bibitem{latentops2023} Liu et al., ``Composable Text Controls in Latent Space with ODEs,'' ACL 2023.
\bibitem{diffuser2023} Reid et al., ``DiffusER: Diffusion via Edit-Based Reconstruction,'' ICLR 2023.
\bibitem{simpleguidance2025} Nisonoff et al., ``Simple Guidance Mechanisms for Discrete Diffusion Models,'' ICLR 2025.
\bibitem{remdm2025} Arriola et al., ``Scaling up Test-Time Compute with Latent Reasoning: A Recurrent Depth Approach,'' arXiv 2025.
\bibitem{drakes2025} Wang et al., ``Fine-Tuning Discrete Diffusion Models via Reward Optimization with Applications to DNA and Protein Design,'' ICLR 2025.
\bibitem{ddpo2024} Black et al., ``Training Diffusion Models with Reinforcement Learning,'' ICLR 2024.
\bibitem{universalguidance2024} Bansal et al., ``Universal Guidance for Diffusion Models,'' ICLR 2024.
\bibitem{sd3_2024} Esser et al., ``Scaling Rectified Flow Transformers for High-Resolution Image Synthesis,'' ICML 2024.
\bibitem{transfusion2025} Zhou et al., ``Transfusion: Predict the Next Token and Diffuse Images with One Multi-Modal Model,'' ICLR 2025.
\bibitem{duo2025} Sahoo et al., ``The Diffusion Duality,'' ICML 2025.
\bibitem{tabdiff2025} Xu et al., ``Mixed-Type Tabular Data Synthesis with Score-based Diffusion in Latent Space,'' ICLR 2025.
\bibitem{ccdd2025} [Authors], ``CCDD: Coupled Continuous-Discrete Diffusion,'' ICLR 2026 submission (rejected).
\bibitem{sdedit2022} Meng et al., ``SDEdit: Guided Image Synthesis and Editing with Stochastic Differential Equations,'' ICLR 2022.
\bibitem{dhariwal2021} Dhariwal and Nichol, ``Diffusion Models Beat GANs on Image Synthesis,'' NeurIPS 2021.
\bibitem{ho2022cfg} Ho and Salimans, ``Classifier-Free Diffusion Guidance,'' NeurIPS Workshop 2022.
\bibitem{cosmos2025} Meshchaninov et al., ``Cosmos: A Smooth Latent Space for Text Diffusion,'' NeurIPS 2025.

\end{thebibliography}

\end{document}
